\chapter{Gram-Schmidt and QR}

There is a recurring theme in this book: \emph{perpendicular is easier}. Whenever
two directions are at right angles, the bookkeeping collapses. Lengths add by
Pythagoras instead of by the law of cosines. Projections decouple, one coordinate
at a time, instead of tangling together in a linear system. Coordinates can be read
off by a single dot product instead of solved for.

So here is a natural wish. Suppose you are handed a few vectors that span a subspace
-- the columns of a matrix, say -- but they point in awkward, oblique directions,
leaning on one another. Could you \emph{straighten them out}? Could you replace your
crooked, redundant frame with a clean orthonormal one that spans exactly the same
space, built directly from the vectors you started with?

You can. The procedure is \textbf{Gram--Schmidt}, and it is nothing more than
``project away what you already have, keep what is new.'' When you record the
arithmetic Gram--Schmidt performs, that record \emph{is} a factorization of your
original matrix: $\mat{A} = \mat{Q}\mat{R}$, with $\mat{Q}$ orthonormal and $\mat{R}$
upper triangular. This is the \textbf{QR factorization}, one of the workhorses of
numerical linear algebra and the engine inside every stable least-squares solver.

This chapter builds that whole story from a picture.

\section{The wish: straighten a crooked basis}

Picture two vectors in the plane, $\vec{a}_1$ and $\vec{a}_2$, that are independent
but \emph{not} perpendicular -- they lean toward each other at maybe forty degrees
(Figure~\ref{ch20:crooked}). They span the whole plane $\R^2$, so as a basis they
are perfectly fine. But they are an \emph{annoying} basis. To find the coordinates of
some vector $\vec{b}$ in this basis, you must solve a $2\times 2$ system, because the
two directions interfere: pushing along $\vec{a}_1$ also moves you partly along
$\vec{a}_2$.

\begin{figure}[H]
\centering
\begin{tikzpicture}[scale=1.3,>=Stealth]
  \draw[->,gray!50] (-0.4,0)--(3,0);
  \draw[->,gray!50] (0,-0.4)--(0,2.6);
  \draw[->,very thick,vecblue] (0,0)--(2.6,0.6) node[right]{$\vec{a}_1$};
  \draw[->,very thick,vecred]  (0,0)--(1.6,2.0) node[above]{$\vec{a}_2$};
  \draw[->,thick,structurepurple,dashed] (0,0)--(0.78,0.18);
  \node[structurepurple] at (1.15,-0.18){\small leaning};
  \draw (0.45,0.10) arc (12:51:0.45);
  \node at (0.75,0.42){\small $\theta$};
\end{tikzpicture}
\caption{A crooked basis. The vectors $\vec{a}_1,\vec{a}_2$ span the plane but lean on
each other at angle $\theta$. We want a perpendicular frame spanning the same plane.}
\label{ch20:crooked}
\end{figure}

\begin{intuitionbox}
An orthonormal basis is a basis where every vector has length $1$ and every pair is
perpendicular. With one, geometry becomes effortless: the coordinate of $\vec{b}$
along direction $\vec{q}_i$ is just the dot product $\inner{\vec{q}_i}{\vec{b}}$, no
system to solve. Gram--Schmidt is the machine that turns \emph{any} basis into an
orthonormal one, building each new direction out of the old ones.
\end{intuitionbox}

Recall from the chapter on projections that the projection of a vector $\vec{a}$ onto
the line through a vector $\vec{q}$ is
\[
  \proj_{\vec{q}}(\vec{a}) \;=\; \frac{\inner{\vec{q}}{\vec{a}}}{\inner{\vec{q}}{\vec{q}}}\,\vec{q}.
\]
If $\vec{q}$ is already a unit vector, the denominator is $1$ and this simplifies to
$\inner{\vec{q}}{\vec{a}}\,\vec{q}$. That single formula is the only tool we need.

\section{Gram-Schmidt for two vectors}

Start with independent vectors $\vec{a}_1$ and $\vec{a}_2$. We want an orthonormal
pair $\vec{q}_1,\vec{q}_2$ spanning the same space.

\textbf{Step 1: the first direction is free.} There is nothing yet to be
perpendicular to, so just keep $\vec{a}_1$'s direction and scale it to length one:
\[
  \vec{q}_1 \;=\; \frac{\vec{a}_1}{\norm{\vec{a}_1}}.
\]

\textbf{Step 2: remove the part of $\vec{a}_2$ that lies along $\vec{q}_1$.}
The vector $\vec{a}_2$ has a shadow on the line of $\vec{q}_1$ -- the component that
$\vec{q}_1$ already covers. Subtract it off. What remains is exactly the part of
$\vec{a}_2$ that points in a genuinely \emph{new} direction, perpendicular to
$\vec{q}_1$:
\[
  \vec{v}_2 \;=\; \vec{a}_2 - \proj_{\vec{q}_1}(\vec{a}_2)
            \;=\; \vec{a}_2 - \inner{\vec{q}_1}{\vec{a}_2}\,\vec{q}_1.
\]
Then normalize:
\[
  \vec{q}_2 \;=\; \frac{\vec{v}_2}{\norm{\vec{v}_2}}.
\]

\begin{figure}[H]
\centering
\begin{tikzpicture}[scale=1.5,>=Stealth]
  \draw[->,gray!40] (-0.3,0)--(3,0);
  \draw[->,gray!40] (0,-0.3)--(0,2.3);
  % q1 along x-ish
  \draw[->,very thick,vecblue] (0,0)--(2.4,0) node[below right]{$\vec{q}_1$};
  % a2 leaning
  \draw[->,very thick,vecred]  (0,0)--(1.5,1.7) node[above]{$\vec{a}_2$};
  % projection (shadow)
  \draw[->,thick,structurepurple] (0,0)--(1.5,0)
       node[midway,below]{\small $\proj_{\vec{q}_1}(\vec{a}_2)$};
  % the perpendicular leftover v2
  \draw[->,very thick,targetgreen] (1.5,0)--(1.5,1.7);
  \node[targetgreen] at (1.92,0.85){$\vec{v}_2$};
  % the drop
  \draw[gray,dashed] (1.5,1.7)--(1.5,0);
  \draw (1.35,0)--(1.35,0.15)--(1.5,0.15);
\end{tikzpicture}
\caption{Gram--Schmidt on two vectors. Subtract the shadow of $\vec{a}_2$ on
$\vec{q}_1$; the leftover $\vec{v}_2$ is automatically perpendicular to $\vec{q}_1$.
Normalize each to get $\vec{q}_1,\vec{q}_2$.}
\label{ch20:gs2}
\end{figure}

Why is $\vec{v}_2$ guaranteed perpendicular to $\vec{q}_1$? Take the dot product and
watch it cancel. Using $\inner{\vec{q}_1}{\vec{q}_1}=1$,
\[
  \inner{\vec{q}_1}{\vec{v}_2}
   = \inner{\vec{q}_1}{\vec{a}_2} - \inner{\vec{q}_1}{\vec{a}_2}\,\inner{\vec{q}_1}{\vec{q}_1}
   = \inner{\vec{q}_1}{\vec{a}_2} - \inner{\vec{q}_1}{\vec{a}_2}
   = 0.
\]
That cancellation is the whole idea: we subtracted off \emph{exactly} the amount that
overlapped $\vec{q}_1$, so nothing of $\vec{q}_1$ is left in $\vec{v}_2$.

\begin{labbox}
Let $\vec{a}_1=\colvec{3\\0}$ and $\vec{a}_2=\colvec{2\\2}$.
\begin{itemize}
  \item $\norm{\vec{a}_1}=3$, so $\vec{q}_1=\colvec{1\\0}$.
  \item $\inner{\vec{q}_1}{\vec{a}_2}=2$, so the shadow is $2\,\vec{q}_1=\colvec{2\\0}$.
  \item $\vec{v}_2=\colvec{2\\2}-\colvec{2\\0}=\colvec{0\\2}$, and
        $\vec{q}_2=\colvec{0\\1}$.
\end{itemize}
Check: $\inner{\vec{q}_1}{\vec{q}_2}=0$, both unit length. We recovered the standard
axes -- which makes sense, since $\vec{a}_1$ was already on the $x$-axis. Try it again
with $\vec{a}_1=\colvec{1\\1}$ and see the frame come out tilted by $45^\circ$.
\end{labbox}

\section{Gram-Schmidt in general}

The pattern now repeats. When we reach $\vec{a}_k$, we already have an orthonormal set
$\vec{q}_1,\dots,\vec{q}_{k-1}$. We subtract off \emph{all} the shadows at once -- the
component of $\vec{a}_k$ along each $\vec{q}_i$ we have built so far -- and whatever
survives is perpendicular to every one of them.

\begin{formalbox}
\textbf{The Gram--Schmidt process.} Given independent vectors
$\vec{a}_1,\dots,\vec{a}_n$, define for $k=1,2,\dots,n$:
\[
  \vec{v}_k \;=\; \vec{a}_k \;-\; \sum_{i=1}^{k-1} \inner{\vec{q}_i}{\vec{a}_k}\,\vec{q}_i,
  \qquad
  \vec{q}_k \;=\; \frac{\vec{v}_k}{\norm{\vec{v}_k}}.
\]
Then $\vec{q}_1,\dots,\vec{q}_n$ is an orthonormal set with
$\spn\{\vec{q}_1,\dots,\vec{q}_k\}=\spn\{\vec{a}_1,\dots,\vec{a}_k\}$ for every $k$.
\end{formalbox}

The summation is the key line. It says: \emph{look at how much of $\vec{a}_k$ each
finished direction already explains, and subtract every one of those pieces.} The
remainder $\vec{v}_k$ is the part of $\vec{a}_k$ that none of the earlier directions
could reach.

\begin{focusbox}
Hold onto three words: \textbf{project, subtract, normalize}. Project $\vec{a}_k$ onto
the space you have already built, subtract that projection to expose the new
direction, normalize to unit length. Every line of Gram--Schmidt is one of those three
moves.
\end{focusbox}

Let us prove the two claims -- orthogonality and equal spans -- because both follow in
a few honest lines and they explain \emph{why} the recipe is correct rather than just
asserting it.

\begin{proposition}[Orthogonality]
\label{ch20:orth-prop}
For each $k$, the vector $\vec{v}_k$ defined above is orthogonal to every
$\vec{q}_j$ with $j<k$.
\end{proposition}

\begin{proof}
Fix $j<k$ and dot $\vec{q}_j$ against $\vec{v}_k$:
\[
  \inner{\vec{q}_j}{\vec{v}_k}
   = \inner{\vec{q}_j}{\vec{a}_k}
     - \sum_{i=1}^{k-1}\inner{\vec{q}_i}{\vec{a}_k}\,\inner{\vec{q}_j}{\vec{q}_i}.
\]
By induction the earlier vectors $\vec{q}_1,\dots,\vec{q}_{k-1}$ are already
orthonormal, so $\inner{\vec{q}_j}{\vec{q}_i}=0$ except when $i=j$, where it equals
$1$. The sum therefore collapses to its single $i=j$ term:
\[
  \inner{\vec{q}_j}{\vec{v}_k}
   = \inner{\vec{q}_j}{\vec{a}_k} - \inner{\vec{q}_j}{\vec{a}_k}\cdot 1 = 0. \qedhere
\]
\end{proof}

\begin{proposition}[Equal spans]
\label{ch20:span-prop}
For each $k$, $\spn\{\vec{q}_1,\dots,\vec{q}_k\}=\spn\{\vec{a}_1,\dots,\vec{a}_k\}$.
\end{proposition}

\begin{proof}
Each $\vec{q}_k$ is built only from $\vec{a}_k$ and the earlier
$\vec{q}_1,\dots,\vec{q}_{k-1}$ -- which in turn came from $\vec{a}_1,\dots,\vec{a}_{k-1}$
-- so every $\vec{q}_k$ is a combination of $\vec{a}_1,\dots,\vec{a}_k$. Conversely,
rearranging the definition gives
\[
  \vec{a}_k \;=\; \norm{\vec{v}_k}\,\vec{q}_k \;+\; \sum_{i=1}^{k-1}\inner{\vec{q}_i}{\vec{a}_k}\,\vec{q}_i,
\]
so every $\vec{a}_k$ is a combination of $\vec{q}_1,\dots,\vec{q}_k$. Each spanning set
sits inside the other; the spans are equal. (Note $\norm{\vec{v}_k}\neq 0$ because the
$\vec{a}_k$ are independent, so no projection ever lands exactly on $\vec{a}_k$.)
\end{proof}

\begin{trapbox}
\textbf{Subtract from the original, not from the running remainder} -- and use the
\emph{normalized} $\vec{q}_i$. The clean formula $\inner{\vec{q}_i}{\vec{a}_k}$ assumes
$\vec{q}_i$ has unit length. If you skip normalizing and keep raw $\vec{v}_i$ instead,
you must restore the denominator $\inner{\vec{v}_i}{\vec{v}_i}$ in every projection.
Both versions are correct; mixing them is the classic bug.
\end{trapbox}

\begin{walkbox}
You have built $\vec{q}_1,\dots,\vec{q}_{k-1}$ and you are about to process
$\vec{a}_k$. Before computing anything: what does the leftover $\vec{v}_k$ \emph{mean}
geometrically? (Hint: stack the $\vec{q}_i$ into a subspace and think about the
perpendicular from $\vec{a}_k$ to that subspace -- it is the same projection picture
from the projections chapter, just onto a span instead of a single line.)
\end{walkbox}

Figure~\ref{ch20:gsplane} makes that perpendicular concrete. The shaded plane is the
span of the directions already finished; the sum of shadows
$\sum_i \inner{\vec{q}_i}{\vec{a}_k}\,\vec{q}_i$ is the foot of the perpendicular from
$\vec{a}_k$, and $\vec{v}_k$ is the perpendicular itself -- the part of $\vec{a}_k$ that
sticks out of the old subspace.

\begin{figure}[H]
\centering
\begin{tikzpicture}[scale=1.0,>=Stealth,line join=round]
  \tikzset{x={(1cm,0cm)},y={(0.5cm,0.32cm)},z={(0cm,1cm)}}
  % the already-built subspace as a shaded parallelogram in the x-y "floor"
  \fill[sourceblue!12] (-0.4,-0.4,0)--(3.0,-0.4,0)--(3.0,2.2,0)--(-0.4,2.2,0)--cycle;
  \node[sourceblue] at (2.6,2.0,0){\small $\spn\{\vec{q}_1,\dots,\vec{q}_{k-1}\}$};
  % two finished directions lying in the floor
  \draw[->,very thick,vecblue] (0,0,0)--(1.6,0,0) node[anchor=north]{$\vec{q}_1$};
  \draw[->,very thick,vecblue] (0,0,0)--(0,1.6,0) node[anchor=west]{$\vec{q}_{k-1}$};
  % a_k sticking up out of the plane
  \draw[->,very thick,vecred] (0,0,0)--(1.3,1.1,1.5) node[anchor=south west]{$\vec{a}_k$};
  % its foot (the sum of shadows) in the plane
  \draw[->,thick,structurepurple] (0,0,0)--(1.3,1.1,0);
  \node[structurepurple] at (1.5,0.55,0){\small $\sum_i\inner{\vec{q}_i}{\vec{a}_k}\vec{q}_i$};
  % the perpendicular leftover v_k
  \draw[->,very thick,targetgreen] (1.3,1.1,0)--(1.3,1.1,1.5);
  \node[targetgreen] at (1.75,1.1,0.8){$\vec{v}_k$};
  % drop line + right-angle hint
  \draw[gray,dashed] (1.3,1.1,1.5)--(1.3,1.1,0);
  \draw (1.15,0.95,0)--(1.15,0.95,0.18)--(1.3,1.1,0.18);
\end{tikzpicture}
\caption{Gram--Schmidt in general is projection onto a \emph{plane} (or higher
subspace), not just a line. The foot of the perpendicular is the sum of the shadows on
the finished directions; the leftover $\vec{v}_k$ stands straight up out of
$\spn\{\vec{q}_1,\dots,\vec{q}_{k-1}\}$, then gets normalized to $\vec{q}_k$.}
\label{ch20:gsplane}
\end{figure}

\section{A fully worked three-vector example}

Let us run the whole machine once, slowly, on three vectors in $\R^3$. Take
\[
  \vec{a}_1=\colvec{1\\1\\0},\qquad
  \vec{a}_2=\colvec{1\\0\\1},\qquad
  \vec{a}_3=\colvec{0\\1\\1}.
\]

\textbf{First direction.} $\norm{\vec{a}_1}=\sqrt{1^2+1^2+0^2}=\sqrt2$, so
\[
  \vec{q}_1 = \frac{1}{\sqrt2}\colvec{1\\1\\0}.
\]

\textbf{Second direction.} First the overlap with $\vec{q}_1$:
\[
  \inner{\vec{q}_1}{\vec{a}_2}
   = \frac{1}{\sqrt2}(1\cdot 1 + 1\cdot 0 + 0\cdot 1) = \frac{1}{\sqrt2}.
\]
Subtract that shadow from $\vec{a}_2$:
\[
  \vec{v}_2
   = \colvec{1\\0\\1} - \frac{1}{\sqrt2}\cdot\frac{1}{\sqrt2}\colvec{1\\1\\0}
   = \colvec{1\\0\\1} - \frac12\colvec{1\\1\\0}
   = \colvec{\tfrac12\\[2pt]-\tfrac12\\[2pt]1}.
\]
Its length is $\norm{\vec{v}_2}=\sqrt{\tfrac14+\tfrac14+1}=\sqrt{\tfrac32}=\dfrac{\sqrt6}{2}$, so
\[
  \vec{q}_2 = \frac{2}{\sqrt6}\colvec{\tfrac12\\[2pt]-\tfrac12\\[2pt]1}
            = \frac{1}{\sqrt6}\colvec{1\\-1\\2}.
\]
Quick check: $\inner{\vec{q}_1}{\vec{q}_2}=\frac{1}{\sqrt2}\frac{1}{\sqrt6}(1-1+0)=0$.
Perpendicular, as promised.

\textbf{Third direction.} Now subtract \emph{both} shadows from $\vec{a}_3$. Compute
the two overlaps:
\[
  \inner{\vec{q}_1}{\vec{a}_3}=\frac{1}{\sqrt2}(0+1+0)=\frac{1}{\sqrt2},
  \qquad
  \inner{\vec{q}_2}{\vec{a}_3}=\frac{1}{\sqrt6}(0-1+2)=\frac{1}{\sqrt6}.
\]
Then
\[
  \vec{v}_3
   = \colvec{0\\1\\1}
     - \frac{1}{\sqrt2}\cdot\frac{1}{\sqrt2}\colvec{1\\1\\0}
     - \frac{1}{\sqrt6}\cdot\frac{1}{\sqrt6}\colvec{1\\-1\\2}.
\]
The coefficients are $\tfrac12$ and $\tfrac16$, so
\[
  \vec{v}_3
   = \colvec{0\\1\\1}
     - \frac12\colvec{1\\1\\0}
     - \frac16\colvec{1\\-1\\2}
   = \colvec{0-\tfrac12-\tfrac16 \\[2pt] 1-\tfrac12+\tfrac16 \\[2pt] 1-0-\tfrac13}
   = \colvec{-\tfrac23\\[2pt]\tfrac23\\[2pt]\tfrac23}.
\]
Its length is $\norm{\vec{v}_3}=\sqrt{3\cdot\tfrac49}=\sqrt{\tfrac43}=\dfrac{2}{\sqrt3}$, so
\[
  \vec{q}_3 = \frac{\sqrt3}{2}\colvec{-\tfrac23\\[2pt]\tfrac23\\[2pt]\tfrac23}
            = \frac{1}{\sqrt3}\colvec{-1\\1\\1}.
\]
Final sanity check on the third vector against the first two:
\[
  \inner{\vec{q}_1}{\vec{q}_3}=\tfrac{1}{\sqrt6}(-1+1+0)=0,
  \qquad
  \inner{\vec{q}_2}{\vec{q}_3}=\tfrac{1}{\sqrt{18}}(-1-1+2)=0.
\]
Three mutually perpendicular unit vectors, built entirely from the crooked originals.

\begin{figure}[H]
\centering
\begin{tikzpicture}[scale=2.0,>=Stealth,line join=round]
  % simple oblique projection of 3D axes
  \tikzset{x={(1cm,0cm)},y={(0.45cm,0.30cm)},z={(0cm,1cm)}}
  \draw[->,gray!45] (0,0,0)--(1.5,0,0) node[anchor=north]{\small $e_1$};
  \draw[->,gray!45] (0,0,0)--(0,1.5,0) node[anchor=west]{\small $e_2$};
  \draw[->,gray!45] (0,0,0)--(0,0,1.5) node[anchor=south]{\small $e_3$};
  % q1, q2, q3 directions (use the unnormalized-ish directions for shape)
  \draw[->,very thick,vecblue]   (0,0,0)--(0.9,0.9,0)   node[anchor=west]{$\vec{q}_1$};
  \draw[->,very thick,vecred]    (0,0,0)--(0.5,-0.5,1.0) node[anchor=south]{$\vec{q}_2$};
  \draw[->,very thick,targetgreen](0,0,0)--(-0.6,0.6,0.6) node[anchor=east]{$\vec{q}_3$};
  % right-angle ticks (suggestive)
  \node at (0.55,0.05,0.6){\small mutually $\perp$};
\end{tikzpicture}
\caption{The orthonormal frame $\vec{q}_1,\vec{q}_2,\vec{q}_3$ produced from
$\vec{a}_1,\vec{a}_2,\vec{a}_3$. Same span as the originals, but now a perfect
right-angled set.}
\label{ch20:gs3}
\end{figure}

\section{The QR factorization}

Here is the beautiful part. We never threw away the arithmetic. Look again at the
rearranged formula from Proposition~\ref{ch20:span-prop}:
\[
  \vec{a}_k = \inner{\vec{q}_1}{\vec{a}_k}\,\vec{q}_1
            + \dots
            + \inner{\vec{q}_{k-1}}{\vec{a}_k}\,\vec{q}_{k-1}
            + \norm{\vec{v}_k}\,\vec{q}_k.
\]
Every original vector $\vec{a}_k$ is written as a combination of
$\vec{q}_1,\dots,\vec{q}_k$ -- and crucially, \emph{only} those, never the later
$\vec{q}$'s, because $\vec{q}_{k+1},\dots$ did not exist yet when we processed
$\vec{a}_k$. That ``only the earlier $\vec{q}$'s'' fact is exactly what makes a matrix
\emph{upper triangular}.

Stack the $\vec{a}_k$ as the columns of a matrix $\mat{A}$ and the $\vec{q}_k$ as the
columns of $\mat{Q}$. The coefficients above fill a matrix $\mat{R}$, and the relation
becomes a single clean equation:
\[
  \boxed{\;\mat{A} = \mat{Q}\mat{R}\;}
\]
where $\mat{Q}$ has orthonormal columns and $\mat{R}$ is upper triangular with the
entries
\[
  R_{ik} = \inner{\vec{q}_i}{\vec{a}_k}\ \ (i<k),
  \qquad
  R_{kk} = \norm{\vec{v}_k},
  \qquad
  R_{ik}=0\ \ (i>k).
\]

\begin{intuitionbox}
$\mat{A}=\mat{Q}\mat{R}$ separates two jobs that the columns of $\mat{A}$ were doing at
once. $\mat{Q}$ holds the clean \emph{directions} (an orthonormal frame); $\mat{R}$
holds the \emph{recipe} -- how much of each direction each original column needed, and
how long each one was. Reading column $k$ of $\mat{A}=\mat{Q}\mat{R}$ is reading
``$\vec{a}_k$ = (these amounts) of (these clean directions).''
\end{intuitionbox}

The triangularity of $\mat{R}$ is not a coincidence to be memorized -- it is the
direct fingerprint of the \emph{order} in which Gram--Schmidt works. The diagonal
entries $R_{kk}=\norm{\vec{v}_k}$ are positive, which is what records that each new
$\vec{a}_k$ genuinely added a new dimension.

\subsection{QR for our three-vector example}

We already have everything. The dot products we computed are precisely the entries of
$\mat{R}$. With
\[
  \mat{Q}=
  \begin{bmatrix}
    \tfrac{1}{\sqrt2} & \tfrac{1}{\sqrt6} & -\tfrac{1}{\sqrt3}\\[4pt]
    \tfrac{1}{\sqrt2} & -\tfrac{1}{\sqrt6} & \tfrac{1}{\sqrt3}\\[4pt]
    0 & \tfrac{2}{\sqrt6} & \tfrac{1}{\sqrt3}
  \end{bmatrix},
  \qquad
  \mat{R}=
  \begin{bmatrix}
    \norm{\vec{a}_1} & \inner{\vec{q}_1}{\vec{a}_2} & \inner{\vec{q}_1}{\vec{a}_3}\\[4pt]
    0 & \norm{\vec{v}_2} & \inner{\vec{q}_2}{\vec{a}_3}\\[4pt]
    0 & 0 & \norm{\vec{v}_3}
  \end{bmatrix}
  =
  \begin{bmatrix}
    \sqrt2 & \tfrac{1}{\sqrt2} & \tfrac{1}{\sqrt2}\\[4pt]
    0 & \tfrac{\sqrt6}{2} & \tfrac{1}{\sqrt6}\\[4pt]
    0 & 0 & \tfrac{2}{\sqrt3}
  \end{bmatrix}.
\]
Let us verify the first column of $\mat{Q}\mat{R}$ to make sure nothing slipped. Column
$1$ of $\mat{R}$ is $(\sqrt2,0,0)\T$, so
\[
  \mat{Q}\,\colvec{\sqrt2\\0\\0}
  = \sqrt2\,\vec{q}_1
  = \sqrt2\cdot\frac{1}{\sqrt2}\colvec{1\\1\\0}
  = \colvec{1\\1\\0}
  = \vec{a}_1. \checkmark
\]
And column $2$: $\mat{R}$'s second column is
$(\tfrac{1}{\sqrt2},\tfrac{\sqrt6}{2},0)\T$, giving
\[
  \frac{1}{\sqrt2}\vec{q}_1 + \frac{\sqrt6}{2}\vec{q}_2
  = \frac12\colvec{1\\1\\0} + \frac{\sqrt6}{2}\cdot\frac{1}{\sqrt6}\colvec{1\\-1\\2}
  = \frac12\colvec{1\\1\\0} + \frac12\colvec{1\\-1\\2}
  = \colvec{1\\0\\1}
  = \vec{a}_2. \checkmark
\]

\begin{labbox}
Factor $\mat{A}=\begin{bmatrix}1&2\\0&1\\1&0\end{bmatrix}$ by hand. The columns are
$\vec{a}_1=(1,0,1)\T$ and $\vec{a}_2=(2,1,0)\T$.
\begin{itemize}
  \item $\norm{\vec{a}_1}=\sqrt2$, $\vec{q}_1=\tfrac{1}{\sqrt2}(1,0,1)\T$.
  \item $\inner{\vec{q}_1}{\vec{a}_2}=\tfrac{1}{\sqrt2}(2+0+0)=\sqrt2$.
  \item $\vec{v}_2=(2,1,0)\T-\sqrt2\cdot\tfrac{1}{\sqrt2}(1,0,1)\T=(1,1,-1)\T$,
        $\norm{\vec{v}_2}=\sqrt3$, $\vec{q}_2=\tfrac{1}{\sqrt3}(1,1,-1)\T$.
\end{itemize}
Then $\mat{R}=\begin{bmatrix}\sqrt2 & \sqrt2\\ 0 & \sqrt3\end{bmatrix}$. Multiply
$\mat{Q}\mat{R}$ back out and confirm you recover $\mat{A}$.
\end{labbox}

\subsection{Why \texorpdfstring{$\mat{Q}$}{Q} is so pleasant}

A matrix whose columns are orthonormal satisfies the single most useful identity in
this corner of the subject:
\[
  \mat{Q}\T\mat{Q} = \mat{I}.
\]
Entry $(i,j)$ of $\mat{Q}\T\mat{Q}$ is $\inner{\vec{q}_i}{\vec{q}_j}$, which is $1$ when
$i=j$ and $0$ otherwise -- that is the identity matrix exactly. When $\mat{Q}$ is
square it is called \textbf{orthogonal}, and then $\mat{Q}\T=\mat{Q}\inv$: the inverse
is free, just the transpose. Orthogonal matrices preserve lengths and angles
($\norm{\mat{Q}\vec{x}}=\norm{\vec{x}}$), so multiplying by $\mat{Q}$ is a rigid motion
-- a rotation or reflection. This is the deep reason QR is numerically gentle: it never
stretches errors.

\begin{strangbox}
Strang introduces orthonormal bases and Gram--Schmidt right alongside projections and
least squares, and he writes $\mat{A}=\mat{Q}\mat{R}$ in essentially the form above. If
you are reading this \emph{before} Strang, the thing to carry in is the picture of
Figure~\ref{ch20:gs2}: project, subtract, normalize. If you are reading \emph{after}
Strang, the payoff is the next two sections -- why $\mat{Q}\T\mat{Q}=\mat{I}$ turns the
normal equations into a triangular solve, and why the modified version is what
software actually runs.
\end{strangbox}

\section{Why orthonormal columns make least squares easy}

Recall the least-squares problem from the previous chapter: $\mat{A}\vec{x}=\vec{b}$
usually has no solution because $\vec{b}$ is not in the column space, so we settle for
the $\hat{\vec{x}}$ that minimizes $\norm{\mat{A}\vec{x}-\vec{b}}$. The classic route
is the \textbf{normal equations}
\[
  \mat{A}\T\mat{A}\,\hat{\vec{x}} = \mat{A}\T\vec{b}.
\]
This is correct, but $\mat{A}\T\mat{A}$ can be a numerically nasty matrix -- it
\emph{squares} the conditioning of $\mat{A}$, amplifying roundoff. (If $\mat{A}$ has
columns that are nearly parallel, $\mat{A}\T\mat{A}$ is nearly singular, and small
errors in the data blow up.) QR sidesteps the squaring entirely.

Substitute $\mat{A}=\mat{Q}\mat{R}$ into the normal equations and watch the
simplifications cascade. First,
\[
  \mat{A}\T\mat{A} = (\mat{Q}\mat{R})\T(\mat{Q}\mat{R})
                   = \mat{R}\T\mat{Q}\T\mat{Q}\mat{R}
                   = \mat{R}\T\mat{R},
\]
using $\mat{Q}\T\mat{Q}=\mat{I}$. Likewise
$\mat{A}\T\vec{b}=\mat{R}\T\mat{Q}\T\vec{b}$. The normal equations become
$\mat{R}\T\mat{R}\,\hat{\vec{x}}=\mat{R}\T\mat{Q}\T\vec{b}$. Cancel the common
$\mat{R}\T$ (it is invertible, since its diagonal entries $\norm{\vec{v}_k}$ are
nonzero) to land at the strikingly simple
\[
  \boxed{\;\mat{R}\,\hat{\vec{x}} = \mat{Q}\T\vec{b}.\;}
\]

\begin{focusbox}
Least squares via QR is two cheap steps. Compute the vector $\mat{Q}\T\vec{b}$ (just
dot products), then solve $\mat{R}\hat{\vec{x}}=\mat{Q}\T\vec{b}$ by
\emph{back-substitution} -- trivial, because $\mat{R}$ is upper triangular. No
$\mat{A}\T\mat{A}$ to form, no squaring of errors.
\end{focusbox}

There is a clean geometric reading too. The projection of $\vec{b}$ onto the column
space of $\mat{A}$ -- the best approximation $\hat{\vec{b}}$ -- is simply
$\hat{\vec{b}}=\mat{Q}\mat{Q}\T\vec{b}$, where $\mat{Q}\mat{Q}\T$ is the projection
matrix onto $\col(\mat{A})$. The orthonormal columns let you read off the projection as
a sum of independent one-dimensional projections,
$\hat{\vec{b}}=\sum_i \inner{\vec{q}_i}{\vec{b}}\,\vec{q}_i$, each $\vec{q}_i$
contributing its own coordinate with no cross-talk. That decoupling is the entire
advantage of working perpendicular (Figure~\ref{ch20:lsproj}).

\begin{figure}[H]
\centering
\begin{tikzpicture}[scale=1.0,>=Stealth,line join=round]
  \tikzset{x={(1cm,0cm)},y={(0.5cm,0.32cm)},z={(0cm,1cm)}}
  % column space = plane spanned by q1,q2
  \fill[sourceblue!12] (-0.4,-0.4,0)--(3.0,-0.4,0)--(3.0,2.2,0)--(-0.4,2.2,0)--cycle;
  \node[sourceblue] at (2.55,2.0,0){\small $\col(\mat{A})$};
  \draw[->,very thick,vecblue] (0,0,0)--(1.5,0,0) node[anchor=north]{$\vec{q}_1$};
  \draw[->,very thick,vecblue] (0,0,0)--(0,1.5,0) node[anchor=west]{$\vec{q}_2$};
  % data vector b out of the plane
  \draw[->,very thick,vecred] (0,0,0)--(1.2,1.0,1.4) node[anchor=south west]{$\vec{b}$};
  % its projection onto the column space
  \draw[->,thick,structurepurple] (0,0,0)--(1.2,1.0,0);
  \node[structurepurple] at (1.5,0.5,0){\small $\hat{\vec{b}}=\mat{Q}\mat{Q}\T\vec{b}$};
  % the residual (error) perpendicular to the plane
  \draw[->,very thick,targetgreen] (1.2,1.0,0)--(1.2,1.0,1.4);
  \node[targetgreen] at (1.7,1.0,0.75){\small $\vec{b}-\hat{\vec{b}}$};
  \draw[gray,dashed] (1.2,1.0,1.4)--(1.2,1.0,0);
  \draw (1.05,0.85,0)--(1.05,0.85,0.18)--(1.2,1.0,0.18);
\end{tikzpicture}
\caption{Least squares is one projection. The best approximation $\hat{\vec{b}}$ is the
shadow of $\vec{b}$ on the column space, $\hat{\vec{b}}=\mat{Q}\mat{Q}\T\vec{b}$, and the
residual $\vec{b}-\hat{\vec{b}}$ stands perpendicular to that space. Orthonormal columns
let you read $\hat{\vec{b}}$ off coordinate by coordinate, with no cross-talk.}
\label{ch20:lsproj}
\end{figure}

\begin{example}[Least squares by QR]
\label{ch20:ls-example}
Fit a line through the three points $(0,1),(1,2),(2,2)$ -- the same kind of fit from
the regression chapter. We want $c+dx\approx y$, so
\[
  \mat{A}=\begin{bmatrix}1&0\\1&1\\1&2\end{bmatrix},\qquad
  \vec{b}=\colvec{1\\2\\2}.
\]
Gram--Schmidt on the columns gives
$\vec{q}_1=\tfrac{1}{\sqrt3}(1,1,1)\T$ and, after subtracting,
$\vec{v}_2=(0,1,2)\T-1\cdot(1,1,1)\T=(-1,0,1)\T$, so
$\vec{q}_2=\tfrac{1}{\sqrt2}(-1,0,1)\T$. Here
$\inner{\vec{q}_1}{\vec{a}_2}=\tfrac{1}{\sqrt3}(0+1+2)=\sqrt3$ and
$\norm{\vec{v}_2}=\sqrt2$, so
\[
  \mat{R}=\begin{bmatrix}\sqrt3 & \sqrt3\\ 0 & \sqrt2\end{bmatrix},
  \qquad
  \mat{Q}\T\vec{b}=
  \colvec{\tfrac{1}{\sqrt3}(1+2+2)\\[4pt]\tfrac{1}{\sqrt2}(-1+0+2)}
  =\colvec{\tfrac{5}{\sqrt3}\\[4pt]\tfrac{1}{\sqrt2}}.
\]
Back-substitute $\mat{R}\hat{\vec{x}}=\mat{Q}\T\vec{b}$. The bottom row gives
$\sqrt2\,d=\tfrac{1}{\sqrt2}$, so $d=\tfrac12$. The top row gives
$\sqrt3\,c+\sqrt3\,d=\tfrac{5}{\sqrt3}$, i.e. $c+d=\tfrac53$, so
$c=\tfrac53-\tfrac12=\tfrac76$. The best line is $y=\tfrac76+\tfrac12 x$, matching what
the normal equations would give -- but reached by a triangular solve.
\end{example}

\begin{applicationbox}
\textbf{QR inside every regression and every deep-net training loop.} When you call
\texttt{numpy.linalg.lstsq}, fit an ordinary least-squares model in a statistics
package, or solve the linear sub-problem inside a Gauss--Newton or
Levenberg--Marquardt optimizer, the library does \emph{not} form $\mat{A}\T\mat{A}$.
It computes a QR factorization (or the closely related SVD) and back-substitutes,
exactly because squaring the condition number can destroy accuracy on real,
near-collinear data. In large-scale machine learning, QR also underlies stable
orthogonalization steps: ``QR-based whitening,'' re-orthogonalizing the basis vectors
in Krylov methods like \textsc{gmres} and Lanczos (the iterative solvers behind huge
sparse systems), and orthogonal weight initialization, where a network's weight matrix
is set to the $\mat{Q}$ of a random matrix so that signals neither explode nor vanish
as they pass through layers. Whenever ``keep the directions independent and the lengths
honest'' matters, QR is the tool.
\end{applicationbox}

\begin{connectionbox}
\textbf{Gram--Schmidt is not really about $\R^n$ -- it is about inner products.} The
process used \emph{only} dot products, lengths, and projections. Swap the dot product
for any other inner product and the identical recipe orthogonalizes those objects. The
famous case: on the space of polynomials with the inner product
$\inner{f}{g}=\int_{-1}^{1} f(x)g(x)\,dx$, applying Gram--Schmidt to the plain powers
$1,x,x^2,x^3,\dots$ produces the \textbf{Legendre polynomials}, the orthogonal basis at
the heart of Gaussian quadrature and spherical harmonics. Other weightings give
Chebyshev, Hermite, and Laguerre families. So this chapter quietly contains the
construction of every classical orthogonal-polynomial system in analysis -- same three
words: project, subtract, normalize.
\end{connectionbox}

\section{Classical versus modified Gram-Schmidt}

The version we derived is \textbf{classical} Gram--Schmidt: to build $\vec{v}_k$,
project the \emph{original} $\vec{a}_k$ onto each finished $\vec{q}_i$ and subtract all
the shadows from $\vec{a}_k$ at once. Mathematically flawless. Numerically, it has a
weakness: each projection coefficient $\inner{\vec{q}_i}{\vec{a}_k}$ is computed from
the same starting $\vec{a}_k$, and when the vectors are nearly dependent, tiny
roundoff errors in the early $\vec{q}_i$ are not corrected, so the computed $\vec{q}_k$
can drift away from being truly orthogonal to the earlier ones.

The fix is small and clever. \textbf{Modified Gram--Schmidt} subtracts the shadows
\emph{one at a time}, each from the running remainder rather than from the original:

\begin{formalbox}
\textbf{Modified Gram--Schmidt.} Set $\vec{v}\leftarrow\vec{a}_k$. Then for
$i=1,2,\dots,k-1$ in order, update
\[
  \vec{v}\;\leftarrow\;\vec{v}-\inner{\vec{q}_i}{\vec{v}}\,\vec{q}_i,
\]
and finally $\vec{q}_k=\vec{v}/\norm{\vec{v}}$.
\end{formalbox}

In exact arithmetic this gives \emph{exactly} the same answer as the classical
version: subtracting components one at a time, each from the partially-cleaned vector,
removes the same total. The difference appears only in floating point. After we
subtract the $\vec{q}_1$-shadow, the remainder $\vec{v}$ already has its
$\vec{q}_1$-component removed (up to roundoff), so when we next compute
$\inner{\vec{q}_2}{\vec{v}}$ we are dotting against a vector that no longer secretly
carries a $\vec{q}_1$ error. Each step measures the overlap against the freshest,
most-orthogonalized remainder, and the errors stop compounding.

\begin{trapbox}
The classical and modified algorithms differ by \emph{which vector you dot against}:
classical always uses the original $\vec{a}_k$; modified uses the running remainder
$\vec{v}$. On a chalkboard with exact fractions you will get identical numbers and may
wonder why anyone bothers. The bother is real and it is purely numerical -- on a
computer with finite precision, modified Gram--Schmidt keeps $\mat{Q}\T\mat{Q}$ far
closer to $\mat{I}$. If you ever code Gram--Schmidt, code the modified form.
\end{trapbox}

\begin{remark}
Even modified Gram--Schmidt is not the gold standard for the worst-conditioned
problems. Production QR routines (\textsc{lapack}'s \texttt{dgeqrf}) instead build
$\mat{Q}$ from a sequence of \textbf{Householder reflections} -- each one a single
orthogonal matrix that zeros out a whole column below the diagonal at once. Householder
QR is backward stable in the strongest sense and is what truly runs underneath
\texttt{numpy} and \textsc{matlab}. Gram--Schmidt remains the version to \emph{learn},
because it makes the geometry transparent: it is literally projection, repeated. We
meet Householder reflections again in the numerical chapter.
\end{remark}

\begin{aestheticbox}
Step back and notice the unity. \emph{One} idea -- subtract the projection -- gives you
an orthonormal basis (Gram--Schmidt), a matrix factorization ($\mat{A}=\mat{Q}\mat{R}$),
a stable least-squares solver ($\mat{R}\hat{\vec{x}}=\mat{Q}\T\vec{b}$), and, by
changing the inner product, every classical family of orthogonal polynomials. The same
geometric move that straightened two leaning arrows in the plane underwrites the
regression in a spreadsheet, the orthogonal initialization of a neural network, and
Gaussian quadrature in a physics code. ``Perpendicular is easier'' is not a slogan; it
is a construction, and Gram--Schmidt is how you build it.
\end{aestheticbox}

\section{Cost and intuition recap}

One move, repeated: \textbf{project, subtract, normalize}. Here is the whole chapter on
one page.

\begin{itemize}
  \item \textbf{The wish.} A crooked basis tangles every coordinate computation into a
        linear solve. An orthonormal basis untangles it: each coordinate is one dot
        product, $\inner{\vec{q}_i}{\vec{b}}$.
  \item \textbf{The construction.} Gram--Schmidt builds the orthonormal frame one
        direction at a time. For each $\vec{a}_k$, subtract its shadows on the finished
        $\vec{q}_i$ and normalize what is left:
        $\vec{v}_k=\vec{a}_k-\sum_{i<k}\inner{\vec{q}_i}{\vec{a}_k}\vec{q}_i$,
        $\vec{q}_k=\vec{v}_k/\norm{\vec{v}_k}$.
  \item \textbf{The bookkeeping is a factorization.} Recording those dot products and
        lengths gives $\mat{A}=\mat{Q}\mat{R}$: $\mat{Q}$ holds clean directions,
        $\mat{R}$ (upper triangular, positive diagonal) holds the recipe.
  \item \textbf{The payoff.} $\mat{Q}\T\mat{Q}=\mat{I}$ turns the normal equations into
        the triangular solve $\mat{R}\hat{\vec{x}}=\mat{Q}\T\vec{b}$, dodging the
        condition-number squaring of $\mat{A}\T\mat{A}$.
\end{itemize}

\textbf{Cost.} For an $m\times n$ matrix ($m\ge n$), Gram--Schmidt (classical or
modified) costs about $2mn^2$ floating-point operations -- the same leading order as
forming and factoring $\mat{A}\T\mat{A}$, but with far better accuracy because the
errors never get squared. On a chalkboard the bottleneck is instead the square roots and
fractions in the normalizations; choosing friendly starting vectors keeps the arithmetic
clean. \textbf{Stability.} Modified Gram--Schmidt keeps $\mat{Q}\T\mat{Q}$ much closer to
$\mat{I}$ than classical; Householder QR (the production default) is better still. The
geometry, though, is identical across all three: subtract the projection.

\section{Exercises}

\begin{exercise}[Two-vector warmup]
Apply Gram--Schmidt to $\vec{a}_1=\colvec{2\\1}$ and $\vec{a}_2=\colvec{3\\4}$. Find
$\vec{q}_1,\vec{q}_2$ and verify they are orthonormal. Then write down $\mat{R}$ and
confirm $\mat{A}=\mat{Q}\mat{R}$.
\end{exercise}

\begin{exercise}[A clean three-vector run]
Orthonormalize
$\vec{a}_1=\colvec{1\\1\\1}$, $\vec{a}_2=\colvec{1\\1\\0}$, $\vec{a}_3=\colvec{1\\0\\0}$.
(They are independent, so the process will not break down.) Report
$\vec{q}_1,\vec{q}_2,\vec{q}_3$ and the matrix $\mat{R}$.
\end{exercise}

\begin{exercise}[Read off $\mat{R}$ without recomputing]
Suppose you have already produced $\vec{q}_1,\vec{q}_2$ from $\vec{a}_1,\vec{a}_2$ and
found $\norm{\vec{a}_1}=5$, $\inner{\vec{q}_1}{\vec{a}_2}=3$, $\norm{\vec{v}_2}=4$.
Write the full $2\times 2$ matrix $\mat{R}$. What is $\det\mat{R}$, and why does it
equal the area of the parallelogram spanned by $\vec{a}_1,\vec{a}_2$?
\end{exercise}

\begin{exercise}[Least squares by QR]
Fit a line to the points $(1,1),(2,2),(3,2),(4,5)$ using the QR route: build
$\mat{A},\vec{b}$, run Gram--Schmidt to get $\mat{Q},\mat{R}$, then solve
$\mat{R}\hat{\vec{x}}=\mat{Q}\T\vec{b}$ by back-substitution. Compare your slope and
intercept against solving the normal equations directly.
\end{exercise}

\begin{exercise}[Order matters for $\mat{Q}$, not for the span]
Run Gram--Schmidt on $\vec{a}_1=\colvec{1\\0\\0},\vec{a}_2=\colvec{1\\1\\0}$ and then
again on the \emph{reversed} list $\vec{a}_2,\vec{a}_1$. You get different orthonormal
bases. Explain geometrically why the two $\vec{q}_1$'s point differently but the spans
of $\{\vec{q}_1,\vec{q}_2\}$ are identical.
\end{exercise}

\begin{exercise}[Why $\mat{Q}\T\mat{Q}=\mat{I}$ but $\mat{Q}\mat{Q}\T\neq\mat{I}$]
Take the $\mat{Q}$ from Example~\ref{ch20:ls-example} (a $3\times 2$ matrix with
orthonormal columns). Compute both $\mat{Q}\T\mat{Q}$ and $\mat{Q}\mat{Q}\T$. One is
the $2\times 2$ identity; the other is a $3\times 3$ projection matrix. Identify which
subspace $\mat{Q}\mat{Q}\T$ projects onto, and explain why it cannot be the full
identity.
\end{exercise}

\begin{exercise}[Classical vs.\ modified, by hand]
Let $\eps$ be tiny and take
$\vec{a}_1=\colvec{1\\\eps\\0},\vec{a}_2=\colvec{1\\0\\\eps},\vec{a}_3=\colvec{1\\0\\0}$
(nearly dependent). Carry out both classical and modified Gram--Schmidt symbolically,
keeping terms to first order in $\eps$. Show that in exact arithmetic both give the same
$\vec{q}_3$, and describe in words where, with rounding, the classical version would let
a $\vec{q}_1$-error leak into $\vec{q}_3$ while the modified version would not.
\end{exercise}

\begin{exercise}[Orthogonal matrices preserve length -- proof]
Let $\mat{Q}$ be square with $\mat{Q}\T\mat{Q}=\mat{I}$. Prove that
$\norm{\mat{Q}\vec{x}}=\norm{\vec{x}}$ for every $\vec{x}$, and that
$\inner{\mat{Q}\vec{x}}{\mat{Q}\vec{y}}=\inner{\vec{x}}{\vec{y}}$ for all $\vec{x},
\vec{y}$. (Start from $\norm{\mat{Q}\vec{x}}^2=(\mat{Q}\vec{x})\T(\mat{Q}\vec{x})$.)
Conclude that multiplication by $\mat{Q}$ is a rotation or reflection, and connect this
to why QR-based computation does not amplify roundoff.
\end{exercise}

\begin{exercise}[Legendre by hand]
Using the inner product $\inner{f}{g}=\int_{-1}^{1}f(x)g(x)\,dx$, apply Gram--Schmidt
to $1,x,x^2$. Confirm you obtain (up to scaling) $1,\ x,\ x^2-\tfrac13$, the first three
Legendre polynomials, and that they are mutually orthogonal under this integral.
\end{exercise}

\begin{exercise}[A non-square challenge]
Let $\mat{A}$ be $4\times 3$ with independent columns. Explain why the ``thin'' QR has
$\mat{Q}$ of size $4\times 3$ and $\mat{R}$ of size $3\times 3$, and why
$\mat{Q}\T\mat{Q}=\mat{I}_3$ while $\mat{Q}\mat{Q}\T$ is a $4\times 4$ projector of rank
$3$. What does the rank-$3$ projector do to a vector $\vec{b}\in\R^4$, and how does that
relate to least squares?
\end{exercise}
